\chapter{Taller: Arduino}

\section{\obj}
\capacidad
\begin{itemize}
    \item Entender la estructura básica de un sketch de Arduino.
    \item Subir un programa a un Arduino UNO R4 Minima.
    \item Usar entradas y salidas digitales.
    \item Usar la comunicación serial para depuración.
\end{itemize}

\section{\mat}
\begin{itemize}
    \item 1 Arduino UNO R4 Minima.
    \item 1 cable USB-C.
    \item 1 protoboard.
    \item 1 LED.
    \item 1 resistencia de 220 $\Omega$.
    \item 1 pulsador.
    \item Cables de conexi\'on.
\end{itemize}

\section{Metodología}

Este taller tiene una duración de 4 lecciones, repartidas en dos semanas. Los estudiantes deben mostrar durante las clases programadas las actividades propuestas. Deben recabar fotografías, código y resultados experimentales para elaborar las evidencias. Las evidencias se subirán al TecDigital la semana siguiente de finalizadas las actividades.

\section{Práctica en Clase}

\subsection{Actividad 1}

Instale y configure el entorno de desarrollo para trabajar con el Arduino UNO R4 Minima.

\subsubsection{Realice lo siguiente:}
\begin{itemize}
    \item Instale Arduino IDE 2.x desde el sitio oficial de Arduino.
    \item Seleccione la tarjeta \textit{Arduino UNO R4 Minima}.
    \item Seleccione el puerto serial correspondiente.
    \item Abra el ejemplo \textit{Blink}.
    \item Identifique las funciones \texttt{setup()} y \texttt{loop()}.
    \item Modifique el tiempo de parpadeo del LED integrado.
\end{itemize}

\subsubsection{Conteste las preguntas}

\textquestiondown Qué función cumple \texttt{setup()}? \textquestiondown Qué función cumple \texttt{loop()}? \textquestiondown Qué ocurre al cambiar el valor de \texttt{delay()}? \textquestiondown Cuál es la diferencia entre cargar un programa y abrir el monitor serial?

\subsection{Actividad 2}

Implemente un circuito de salida digital con un LED externo y un circuito de entrada digital con un pulsador.

\subsubsection{Realice lo siguiente:}
\begin{itemize}
    \item Conecte un LED externo al pin digital 8 con una resistencia limitadora de 220 $\Omega$.
    \item Programe el LED para encender y apagar de forma periódica.
    \item Conecte un pulsador al pin digital 7.
    \item Configure el pin del pulsador con \texttt{INPUT\_PULLUP}.
    \item Modifique el programa para que el pulsador controle el estado del LED.
\end{itemize}

\subsubsection{Conteste las preguntas}

\textquestiondown Por qué el LED requiere una resistencia en serie? \textquestiondown Qué ventaja tiene usar \texttt{INPUT\_PULLUP}? \textquestiondown Por qué el pulsador se lee en nivel bajo cuando se presiona? \textquestiondown Cómo cambia el comportamiento del sistema cuando se elimina la instrucción \texttt{pinMode()}?

\subsection{Actividad 3}

Implemente un programa que combine entrada digital, salida digital y monitoreo serial.

\subsubsection{Realice lo siguiente:}
\begin{itemize}
    \item Inicialice la comunicación serial con \texttt{Serial.begin(9600)}.
    \item Programe el LED para parpadear lentamente en condiciones normales.
    \item Modifique el programa para que, al presionar el pulsador, el LED parpadee rápidamente.
    \item Muestre en el monitor serial el estado del pulsador y la velocidad de parpadeo seleccionada.
\end{itemize}

\subsubsection{Conteste las preguntas}

\textquestiondown Cuál es la diferencia entre \texttt{Serial.print()} y \texttt{Serial.println()}? \textquestiondown Cómo ayuda el monitor serial al proceso de depuración? \textquestiondown Qué ocurre si la velocidad configurada en el programa no coincide con la del monitor serial? \textquestiondown Cómo se relaciona el tiempo de parpadeo con el comportamiento observado del LED?
